\documentclass[11pt]{article}
\usepackage[margin=1in]{geometry}
\usepackage{amsmath,amssymb,booktabs,longtable,tabularx,array}
\usepackage{enumitem}
\usepackage[T1]{fontenc}
\usepackage[hidelinks]{hyperref}
\setlist[itemize]{leftmargin=*}
\setlist[enumerate]{leftmargin=*}
\newcommand{\code}[1]{\texttt{\detokenize{#1}}}
\newcommand{\st}{s_t}
\newcommand{\snext}{s_{t+1}}
\newcommand{\obs}{o_t}
\newcommand{\sobs}{\sigma_{\mathrm{obs}}}
\newcommand{\ssafe}{s_{\mathrm{safe}}}
\newcommand{\E}{\mathbb{E}}
\newcommand{\Var}{\mathrm{Var}}
\newcommand{\1}{\mathbf{1}}
\sloppy


\title{Real Ecology Data, Actions, and Costs}
\author{Ecology Benchmark Documentation}
\date{July 2026}

\begin{document}
\maketitle

\section{Purpose}

The real-ecology setting is defined by six CSV files in
\code{real_ecology_data}. The loader in \code{realdata.py} validates 9
populations, 11 actions, 45 measured growth-rate rows, and 99
population/action effect rows. The environment reads the precomputed long table,
while the relational core remains the provenance layer.

\section{CSV Inventory}

\begin{center}
\begin{tabularx}{\linewidth}{lll}
\toprule
File & Rows & Role \\
\midrule
\code{species.csv} & 9 & population scale, capacity, and rate caps \\
\code{species_lambda.csv} & 45 & five measured $\lambda$ values per population \\
\code{actions.csv} & 11 & population-independent action menu and cost index \\
\code{action_effects_long.csv} & 99 & precomputed population/action effects \\
\code{cost_sources.csv} & 12 & source keys behind cost scaling \\
\code{cost_anchors_portal.csv} & 12 & portal studies with numeric cost anchors \\
\bottomrule
\end{tabularx}
\end{center}

\section{Population Table}

\begin{center}
\begin{tabularx}{\linewidth}{llrrrr}
\toprule
Common name & Scientific name & $N_0$ & $K_{\mathrm{base}}$ & $K_{\max}$ & Class \\
\midrule
Egyptian vulture & \emph{Neophron percnopterus} & 41 & 325 & 650 & sink \\
Bottlenose dolphin & \emph{Tursiops truncatus} & 25 & 35 & 70 & sink \\
Amur tiger & \emph{Panthera tigris altaica} & 200 & 250 & 500 & recoverable \\
Spotted turtle & \emph{Clemmys guttata} & 31 & 31 & 62 & recoverable \\
Asian elephant & \emph{Elephas maximus} & 60 & 120 & 240 & recoverable \\
Puerto Rican parrot & \emph{Amazona vittata} & 78 & 250 & 500 & recoverable \\
Iberian lynx & \emph{Lynx pardinus} & 89 & 165 & 330 & recoverable \\
Jaguar & \emph{Panthera onca} & 325 & 325 & 650 & recoverable \\
Crab-eating fox & \emph{Cerdocyon thous} & 41 & 41 & 82 & recoverable \\
\bottomrule
\end{tabularx}
\end{center}

A sink is defined by \code{r_max_ricker <= 0}. The two sink populations are
reported separately from recoverable populations in the recent result summary.

\section{Action Table}

\begin{center}
\begin{tabularx}{\linewidth}{lllr}
\toprule
Action & Name & Channel & Cost \\
\midrule
a0 & Do Nothing & none & 0.0000 \\
a1 & Sustainable Harvest & rate & -0.0500 \\
a2 & Aggressive Harvest & rate & -0.1000 \\
a3 & Predator / Disease Control & rate & 0.2500 \\
a4 & Breeding / Recruitment Support & rate & 0.5000 \\
a5 & Moderate Restoration & capacity & 0.1875 \\
a6 & Intensive Restoration & capacity & 0.5000 \\
a7 & Integrated Conservation (light) & rate+capacity & 0.4375 \\
a8 & Adaptive Conservation Trial & rate+capacity & 0.7500 \\
a9 & Flagship Conservation Programme & rate+capacity & 1.0000 \\
a10 & Translocation & state & 0.3125 \\
\bottomrule
\end{tabularx}
\end{center}

Actions a1 and a2 are harvest/revenue actions and therefore have negative cost
indices. Action a10 is the only direct-state action: it adds 10 percent of the
population's $N_0$ before growth. Capacity actions use
$\Delta K=(K_{\mathrm{multiplier}}-1)K_{\mathrm{base}}$ and accumulate through
$\kappa$.

\section{Growth-Rate Conversion}

The relational join is:
\[
  \lambda =
  \mathrm{species\_lambda}[\mathrm{population},\lambda\_\mathrm{id}
    =\mathrm{actions.lambda\_source}].
\]
For Ricker, Allee, and regime families,
\[
  r = \log(\lambda).
\]
For theta-logistic cells,
\[
  r = \lambda - 1.
\]
The action effect table stores both converted columns. The real action table
sets $\rho$ to this action-specific value and clips it to the population/family
caps in \code{species.csv}. The heaviest exploitation action supplies the lower
rate cap; the strongest recovery action supplies the upper cap.

\section{Cost Provenance}

The population data and action mechanics come from the Dominik table described
in \code{real_ecology_data/README.md}. Costs come from a separate cost-function
source grounded in the Conservation Costing Portal and two supplementary
per-unit studies. The resulting \code{cost_step} is a relative index, not a
dollar amount. It preserves ordering and spacing while avoiding a reward scale
dominated by orders-of-magnitude cost differences.

The source keys include White et al. (2022) for cost-reporting framing; Adams et
al. (2012), Green et al. (2012), James et al. (1999), Frazee et al. (2003), and
McCarthy et al. (2012) for portal-grounded land/species-management anchors;
Laycock et al. (2009) for species recovery programmes; Yong et al. (2023) for
threat-abatement costs; Weise et al. (2014) for translocation per-animal costs;
and the California condor programme note for high-cost captive-breeding context.

\section{Data Caveats}

The cost index is relative and coarse. Portal categories are broader than the
11-action menu, so some within-category ordering is a modelling choice. The
growth-rate values are action set-points, not direct estimates of a full
transition model under every structural family. The real-data claim is therefore
about a transparent ecology-informed benchmark table, not a complete empirical
population model for each species.

\section*{Verification}

Verified against: \path{real_ecology_data/README.md}, \path{real_ecology_data/species.csv}, \path{real_ecology_data/species_lambda.csv}, \path{real_ecology_data/actions.csv}, \path{real_ecology_data/action_effects_long.csv}, \path{real_ecology_data/cost_sources.csv}, \path{real_ecology_data/cost_anchors_portal.csv}, \path{src/real_ecology_benchmark/realdata.py}, \path{src/real_ecology_benchmark/actions.py} @ be96c36

\end{document}
